\documentclass[11pt,reqno]{amsart}
\usepackage{amssymb}
\usepackage{booktabs}
\usepackage[expansion=false]{microtype}
\usepackage{xcolor}
\usepackage{xurl}
\usepackage[colorlinks=true,linkcolor=blue!50!black,citecolor=blue!50!black,urlcolor=blue!50!black]{hyperref}

\newtheorem{theorem}{Theorem}[section]
\newtheorem{lemma}[theorem]{Lemma}
\newtheorem{proposition}[theorem]{Proposition}
\newtheorem{corollary}[theorem]{Corollary}
\newtheorem{conjecture}[theorem]{Conjecture}
\theoremstyle{definition}
\newtheorem{definition}[theorem]{Definition}
\theoremstyle{remark}
\newtheorem{remark}[theorem]{Remark}
\newtheorem{example}[theorem]{Example}

\newcommand{\erfc}{\operatorname{erfc}}
\newcommand{\erfi}{\operatorname{erfi}}
\newcommand{\erf}{\operatorname{erf}}
\newcommand{\Ci}{\operatorname{Ci}}
\newcommand{\Ein}{\operatorname{Ein}}
\newcommand{\Eip}{\mathcal{E}}
\newcommand{\Li}{\operatorname{Li}}
\newcommand{\arctanh}{\operatorname{arctanh}}
\newcommand{\dd}{\,\mathrm{d}}
\newcommand{\QQ}{\mathbb{Q}}
\newcommand{\Qbar}{\overline{\mathbb{Q}}}
\newcommand{\ZZ}{\mathbb{Z}}
\newcommand{\CC}{\mathbb{C}}
\newcommand{\RR}{\mathbb{R}}
\newcommand{\Mul}{\mathsf{M}}
\newcommand{\CT}{\operatorname*{CT}}
\newcommand{\Res}{\operatorname*{Res}}

\title[Analytic properties of Dirichlet series]{Analytic properties of Dirichlet series\\ of functions analytic at infinity}
\author{Fredrik Johansson}
\date{Companion note, November 2026}

\raggedbottom
\begin{document}

\begin{abstract}
For a sequence $F$ that agrees with a function analytic at infinity, we collect analytic properties of the Dirichlet
series $\zeta_F(s) = \sum_{k=1}^{\infty}F(k)k^{-s}$: the Mellin part and its recurrence, the absence of linear
difference and differential equations in $s$, the values at the non-positive integers, growth and mean values, a
Hardy--Littlewood approximate functional equation, the number of zeros up to a given height, and numerical
experiments on the zeros of Khinchin's zeta function. This note accompanies \cite{Joh26}, where these functions are
computed for holonomic $F$; nothing here needs holonomy, and none of it is used by the algorithms there.
\end{abstract}

\maketitle

\section{Setting}
Throughout, $F$ is analytic at infinity:
\begin{equation}\label{eq:f}
  F(k) = f(1/k) \quad (k \ge k_0), \qquad f(x) = \sum_{n=0}^{\infty}a_nx^n, \qquad |a_n| \le C R^n,
  \quad R < k_0,
\end{equation}
and $n_0 = \min\{n : a_n \ne 0\}$ (we exclude $f = 0$). We write $\zeta_F(s) = \sum_{k=1}^{\infty}F(k)k^{-s}$ and
$\zeta_F(s;k_0) = \sum_{k=k_0}^{\infty}F(k)k^{-s}$. Most of what follows is a routine consequence of the expansion in
Hurwitz zeta values \cite{FV96}, \cite[Proposition~5.6]{Joh26},
\begin{equation}\label{eq:zetavalue}
  \zeta_F(s;k_0) = \sum_{n=0}^{\infty}a_n\,\zeta(s+n,k_0), \qquad \zeta(w,k_0) = \sum_{k=k_0}^{\infty}k^{-w},
\end{equation}
which converges locally uniformly on $\CC \setminus \{1, 0, -1, \dots\}$: for $\operatorname{Re}w = v > 1$,
$|\zeta(w,k_0)| \le k_0^{-v}(1 + k_0/(v-1))$, so the terms are $O((R/k_0)^n)$ for large $n$. Hence $\zeta_F$ is
meromorphic in $\CC$, with at most simple poles at $s = 1-n$, $n \ge n_0$, of residue $a_n$.

The running example is the Khinchin zeta function
\[
  Z(s) = \sum_{k=1}^{\infty}p_kk^{-s}, \qquad p_k = \log_2\Bigl(1 + \frac1{k(k+2)}\Bigr),
\]
where $p_k$ is the probability that a partial quotient of the continued fraction of a random real number equals $k$
\cite{Khi35,Khi64}. Here
\[
  f(x) = \frac{2\log(1+x) - \log(1+2x)}{\log 2}, \qquad a_n = \frac{(-1)^n(2^n - 2)}{n\log 2},
\]
with $k_0 = 3$ and $n_0 = 2$, and $Z'(0) = -\log K$, where $K$ is Khinchin's constant.

Finally, let $G_s(x) = |x|^{-s}F(|x|)$, with Fourier transform $\widehat G_s$. The identity
$\sum_{k\in\ZZ}G_s(k) = \sum_{m\in\ZZ}\widehat G_s(m)$ is formally Poisson summation, and the theta splitting of
\cite[Theorem~3.2]{Joh26} is a smoothed and rigorous form of it \cite[Remark~3.4]{Joh26}. For $F = 1$,
$\widehat G_s$ is again a power, which is the functional equation of $\zeta$.

\begin{remark}[Mellin--Barnes]\label{rem:mb}
If $F$ is defined on all of $(0,\infty)$, with Mellin transform $\mathcal M F(w) = \int_0^\infty F(x)x^{w-1}\dd x$ in a strip
$c_1 < \operatorname{Re}w < n_0$, then by Mellin inversion
\[
  \zeta_F(s) = \frac1{2\pi i}\int_{(c)}\mathcal MF(w)\,\zeta(s+w)\dd w \qquad (c_1 < c < n_0,\ \operatorname{Re}(s+c) > 1).
\]
Unlike $\zeta_F$, the Mellin transform of a holonomic function satisfies a linear recurrence in $w$ with polynomial
coefficients. Moving the contour to the right picks up the poles of $\mathcal MF$ at $w = n$, which come from the
expansion of $F$ at infinity and have residues $-a_n$: this gives \eqref{eq:zetavalue} (for the tail $k \ge k_0$, where
it converges). Moving it to the left picks up the poles coming from the behaviour of $F$ at $0$, i.e.\ of $f$ at
infinity, and gives a second expansion in values $\zeta(s - w_j)$ and their derivatives, typically divergent but
asymptotic. For Khinchin's $Z$, $p(x) = -\log_2(2x) + \sum_jb_jx^j$ near $x = 0$, and the left-hand expansion involves
$\zeta'(s)$, $\zeta(s)$ and $\sum_jb_j\zeta(s-j)$. With a scale parameter, $\sum_kF(ky)k^{-s}$ has Mellin transform
$\mathcal MF(w)\zeta(s+w)$ in $y$, which gives asymptotic expansions in parameters such as the spacing $\lambda$ of the
vortex sum or the argument $x$ of the Hardy--Littlewood function \cite{FGD95,Kat97}; this complements
the theta splitting of \cite[Theorem~3.2]{Joh26}, which is exact at a fixed parameter.
\end{remark}


\section{The Mellin part}
Let $\mathcal M_F(s;k_0) = \int_{k_0}^\infty F(x)x^{-s}\dd x$, the Mellin transform of $F$ restricted to $[k_0,\infty)$.

\begin{proposition}\label{prop:mellin}
Under \eqref{eq:f}, $\mathcal M_F(s;k_0) = \sum_{n=0}^{\infty} a_n k_0^{1-s-n}/(s+n-1)$, and $\zeta_F(s;k_0) - \mathcal M_F(s;k_0)$ is entire.
If moreover $F$ is holonomic, annihilated by $\sum_{i=-r}^{r} x^ir_i(\vartheta_x)$, $\vartheta_x = x\frac{\mathrm{d}}{\mathrm{d}x}$, then $\mathcal M_F$
satisfies the inhomogeneous linear recurrence
\[
  \sum_{i=-r}^{r} r_i(s-i-1)\,\mathcal M_F(s-i;k_0) = k_0^{1-s}\,P(s),
\]
where $P$ is a polynomial whose coefficients are linear in $F(k_0), F'(k_0), \dots$.
\end{proposition}

\begin{proof}
The first formula is termwise integration. By \eqref{eq:zetavalue}, the difference is
$\sum_{n=0}^{\infty}a_n(\zeta(s+n,k_0) - k_0^{1-s-n}/(s+n-1))$, a locally uniformly convergent series of entire
functions (the classical proof is by Euler--Maclaurin summation, see \cite[Section~5.3]{Joh26}).
For the recurrence, $\int_{k_0}^\infty x^{-s}\vartheta_xF\dd x = -k_0^{1-s}F(k_0) + (s-1)\mathcal M_F(s;k_0)$ and multiplication by
$x^i$ shifts $s$ by $-i$; iterate.
\end{proof}

So all the poles of $\zeta_F$, and all of its ``holonomic'' structure in $s$, sit in the Mellin part. The
remainder is entire but, in contrast to $\mathcal M_F$, not holonomic in $s$:

\begin{proposition}\label{prop:nonholonomic}
Let $F$ be any sequence with $F(k) \ne 0$ for infinitely many $k$ and $\sum_{k=1}^{\infty}F(k)k^{-s}$ convergent somewhere.
Then $\zeta_F$ satisfies no relation $\sum_{i=0}^rp_i(s)\zeta_F(s+i) = q(s)$ and no relation
$\sum_{i=0}^rp_i(s)\zeta_F^{(i)}(s) = q(s)$ with polynomials $p_i$ not all zero and $q$ a finite
linear combination of exponentials $c^{-s}$ with polynomial coefficients.
\end{proposition}

\begin{proof}
For $\operatorname{Re}s$ large the left-hand sides are generalized Dirichlet series
\[
  \sum_{k=1}^{\infty} F(k)\,c_k(s)\,k^{-s}, \quad c_k(s) = \sum_{i=0}^{r}p_i(s)k^{-i}
  \ \text{or}\ c_k(s) = \sum_{i=0}^{r}p_i(s)(-\log k)^i,
\]
polynomial in $s$ of bounded degree. By the uniqueness theorem for such series (divide by the highest power of
$s$ and let $\operatorname{Re}s\to\infty$, inductively), $F(k)c_k(s) = 0$ for every $k$ not among the finitely many
frequencies of $q$. Hence $c_k(s) = 0$ for infinitely many $k$, and as a polynomial in $k^{-1}$ (respectively
$\log k$) with infinitely many roots, $\sum_{i=0}^{r}p_i(s)y^i \equiv 0$, so all $p_i = 0$.
\end{proof}

Whether $\zeta_F$ can satisfy an algebraic differential equation is a different matter: $\zeta$ itself does not
(Hilbert, Ostrowski \cite{Ost20}), and we expect the same for every $\zeta_F$ that is not a finite combination of
exponentials, but we have no proof.

\section{Values at non-positive integers}
Let $B_n(x)$ denote the Bernoulli polynomials and $\psi$ the digamma function.

\begin{proposition}\label{prop:negative}
Under \eqref{eq:f}, for every integer $m \ge 0$,
\begin{multline*}
  \CT_{s=-m}\zeta_F(s;k_0) = -\sum_{n=0}^{m}a_n\frac{B_{m-n+1}(k_0)}{m-n+1} - a_{m+1}\psi(k_0) \\
  + \sum_{k=k_0}^{\infty}k^m\Bigl(f(1/k) - \sum_{n=0}^{m+1}a_nk^{-n}\Bigr),
\end{multline*}
where $\CT$ is the constant term of the Laurent expansion (the value, if $a_{m+1} = 0$).
\end{proposition}

\begin{proof}
In \eqref{eq:zetavalue}, $\zeta(-j, k_0) = -B_{j+1}(k_0)/(j+1)$ for $j \ge 0$,
$\zeta(1+\varepsilon,k_0) = 1/\varepsilon - \psi(k_0) + O(\varepsilon)$, and the terms with $n \ge m+2$ sum to the
absolutely convergent last sum.
\end{proof}

So the zeta-regularized value of a divergent sum $\sum_{k=1}^{\infty} k^mF(k)$ is a convergent sum of the same type plus
rational (if the $a_n$ are rational) Bernoulli corrections. For Khinchin's $Z$, Proposition~\ref{prop:negative}
at $m = 1$ gives $\CT_{s=-1}Z(s) = -1/\log2$, so that $\log2\,(K_p)^p = (1-p)^{-1} - 1 + O(1-p)$ for the H\"older means $K_p = Z(-p)^{1/p}$ as $p \to 1^-$ \cite[Proposition~6.1]{Joh26}.

\section{Growth and mean values}
Let $\mu_F(\sigma) = \inf\{\mu : \zeta_F(\sigma+it) \ll |t|^\mu \text{ for } |t| \ge 1\}$, and let $\mu = \mu_\zeta$.

\begin{proposition}\label{prop:growth}
Under \eqref{eq:f}, $\mu_F(\sigma) = \mu(\sigma + n_0)$ for every real $\sigma$.
\end{proposition}

\begin{proof}
Fix $\sigma$ and choose $N$ with $1 - N < \sigma - 1$. By \eqref{eq:zetavalue},
\[
  \zeta_F(s) = \sum_{k<k_0}F(k)k^{-s} + \sum_{n=n_0}^{N-1}a_n\zeta(s+n,k_0) + R_N(s),
\]
where $R_N(s) = \sum_{k=k_0}^{\infty}k^{-s}\bigl(f(1/k) - \sum_{n=0}^{N-1}a_nk^{-n}\bigr)$ is bounded on
$\operatorname{Re}s = \sigma$, as are the finite sums, and $\zeta(w,k_0)$ differs from $\zeta(w)$ by a
bounded finite sum. Since $\mu$ is nonincreasing, $\mu_F(\sigma) \le \mu(\sigma+n_0)$. Conversely
$a_{n_0}\zeta(s+n_0) = \zeta_F(s) - \sum_{n=n_0+1}^{N-1}a_n\zeta(s+n) + O(1)$. If $\mu(\sigma+n_0) = 0$ there is nothing to
prove ($\mu_F \ge 0$ everywhere: $\mu_F$ is convex, and $\mu_F = 0$ far to the right, where $\zeta_F$ is an
absolutely convergent Dirichlet series that is not $o(1)$). If $\mu(\sigma+n_0) > 0$, then
$\mu(\sigma+n_0+1) < \mu(\sigma+n_0)$: with $x = \sigma+n_0$, for $x \le 0$ because $\mu(x) = \frac12 - x$ while
$\mu(x+1) \le \max(-\frac x2, -\frac12 - x)$, and for $x > 0$ because $\mu(x+1) = 0$. Hence
$\mu(\sigma+n_0) \le \max(\mu_F(\sigma), \mu(\sigma+n_0+1))$ forces $\mu(\sigma+n_0) \le \mu_F(\sigma)$.
\end{proof}

\begin{corollary}\label{cor:lindelof}
\begin{enumerate}
\item The Lindel\"of hypothesis holds if and only if $\zeta_F(\frac12 - n_0 + it) \ll |t|^\varepsilon$ for every
$\varepsilon > 0$.
\item $\displaystyle\int_1^T\bigl|\zeta_F(\tfrac12 - n_0 + it)\bigr|^2\dd t = |a_{n_0}|^2\,T\log T + O\bigl(T(\log T)^{1/2}\bigr)$.
\end{enumerate}
\end{corollary}

\begin{proof}
On the line $\operatorname{Re}s = \frac12 - n_0$, $\zeta_F(s) = a_{n_0}\zeta(s+n_0) + O(1)$, since
$\zeta(s+n)$ for $n > n_0$ is evaluated at real part at least $\frac32$. (1) is immediate, and (2) follows from
$\int_1^T|\zeta(\frac12+it)|^2\dd t \sim T\log T$ and the Cauchy--Schwarz inequality for the cross term.
\end{proof}

\section{Large imaginary part}\label{sec:larget}

Throughout this section $F$ is analytic at infinity \eqref{eq:f} (no differential equation is needed),
$s = \sigma + it$ with $t > 0$, and
$\chi(s) = 2^s\pi^{s-1}\sin(\pi s/2)\Gamma(1-s)$ is the factor in $\zeta(s) = \chi(s)\zeta(1-s)$.

\subsection{A Hardy--Littlewood approximate functional equation}

\begin{proposition}\label{prop:hl}
Let $0 < \sigma + n_0 < 1$ be fixed, and let $x, y \ge C$ with $2\pi xy = t$, $x \ge k_0$ and $y \le x$. Then
\begin{multline*}
  \zeta_F(s) = \sum_{1\le k\le x}F(k)k^{-s} + \chi(s)\sum_{1\le m\le y}F\Bigl(\frac{t}{2\pi m}\Bigr)m^{s-1} \\
  + O\bigl(x^{-\sigma-n_0} + t^{1/2-\sigma-n_0}y^{\sigma+n_0-1}\bigr).
\end{multline*}
\end{proposition}

\begin{proof}
Write $F(k) = a_{n_0}k^{-n_0} + r(k)$ for $k \ge k_0$, with $r(k) = O(k^{-n_0-1})$. Since $\sigma + n_0 > 0$, the series
$\sum_kr(k)k^{-s}$ converges absolutely, and its tail beyond $x$ is $O(x^{-\sigma-n_0})$. Hence
\[
  \zeta_F(s) = \sum_{k<k_0}F(k)k^{-s} + a_{n_0}\zeta(s+n_0,k_0) + \sum_{k_0\le k\le x}r(k)k^{-s} + O(x^{-\sigma-n_0}).
\]
By the approximate functional equation of Hardy and Littlewood \cite{HL23}, \cite[Theorem~4.13]{Tit86}, applied at
$s + n_0$, whose real part lies in $(0,1)$,
\begin{multline*}
  \zeta(s+n_0) = \sum_{1\le k\le x}k^{-s-n_0} + \chi(s+n_0)\sum_{1\le m\le y}m^{s+n_0-1} \\
  + O\bigl(x^{-\sigma-n_0} + t^{1/2-\sigma-n_0}y^{\sigma+n_0-1}\bigr),
\end{multline*}
and $\zeta(s+n_0,k_0) = \zeta(s+n_0) - \sum_{k=1}^{k_0-1}k^{-s-n_0}$. Collecting the direct terms gives
$\sum_{k\le x}F(k)k^{-s}$. For the dual terms, Stirling's formula gives
$\chi(s) = (2\pi/t)^{s-1/2}e^{i(t+\pi/4)}(1 + O(1/t))$ \cite[(4.12.3)]{Tit86}, so
$\chi(s+n_0) = \chi(s)(2\pi/t)^{n_0}(1 + O(1/t))$, and for $m \le y$, $X_m = t/(2\pi m) \ge x \ge k_0$ and
\[
  F(X_m) = a_{n_0}\Bigl(\frac{2\pi m}{t}\Bigr)^{n_0} + O\Bigl(\Bigl(\frac{2\pi m}{t}\Bigr)^{n_0+1}\Bigr).
\]
Hence
\[
  a_{n_0}\chi(s+n_0)m^{s+n_0-1} = \chi(s)F(X_m)m^{s-1}
  + O\Bigl(t^{1/2-\sigma}m^{\sigma-1}\Bigl(\frac mt\Bigr)^{n_0}\Bigl(\frac1t + \frac mt\Bigr)\Bigr),
\]
and summing over $m \le y$ the error is $O(t^{1/2-\sigma-n_0}y^{\sigma+n_0-1}\cdot y^2/t)$, which is within the claimed
bound since $y^2/t = y/(2\pi x) \le 1$.
\end{proof}

The proof uses only the leading term of $F$ at infinity, and the dual coefficients could equally be taken as
$a_{n_0}(2\pi m/t)^{n_0}$; with $F(t/2\pi m)$ the formula is natural (the amplitude at the stationary point of the phase
$-t\log x + 2\pi mx$, as in the van der Corput transformation \cite[Lemma~4.10]{Tit86}), it makes the error terms of
the two approximate functional equations in Table~\ref{tab:hl} comparable, and we expect it to hold without the
restriction $y \le x$, with an error $O(x^{-\sigma-n_0} + t^{1/2-\sigma-n_0}y^{\sigma+n_0-1})$.

For $F(k) \sim a_{n_0}k^{-n_0}$ the dual coefficients are $F(t/2\pi m) \approx a_{n_0}(2\pi m/t)^{n_0}$, and since
$\chi(s+n_0) \approx (t/2\pi)^{-n_0}\chi(s)$ the formula reduces to the classical one for $a_{n_0}\zeta(s+n_0)$ to
leading order; the higher coefficients of $f$ enter through the values $F(t/2\pi m)$ at the stationary points.
Table~\ref{tab:hl} compares the error of Proposition~\ref{prop:hl} for Khinchin's $Z$ (with $x = y$) with the
error of the classical formula for $\zeta(s+2)$; without the dual sum the error would be of the size of the
dual sum itself.

\begin{table}[t]
\centering\small
\caption{Proposition~\ref{prop:hl} for $F(k) = p_k\log 2$ ($n_0 = 2$) with $x = y = \sqrt{t/2\pi}$: absolute values of
the function, of the dual sum, and of the error, compared with the error of the classical approximate functional
equation for $\zeta(s+2)$ with the same $x, y$.}
\label{tab:hl}
\begin{tabular}{rrrrrr}
\toprule
$t$ & $\sigma$ & $|Z(s)\log 2|$ & dual sum & error, $Z$ & error, $\zeta(s+2)$ \\
\midrule
$1\,000$ & $-1.5$ & $1.022$ & $0.657$ & $0.098$ & $0.114$ \\
$1\,000$ & $-1.2$ & $0.355$ & $0.187$ & $0.046$ & $0.053$ \\
$4\,000$ & $-1.5$ & $0.341$ & $1.338$ & $0.097$ & $0.104$ \\
$4\,000$ & $-1.2$ & $0.131$ & $0.409$ & $0.037$ & $0.039$ \\
$16\,000$ & $-1.5$ & $0.459$ & $0.547$ & $0.052$ & $0.054$ \\
$16\,000$ & $-1.2$ & $0.324$ & $0.137$ & $0.016$ & $0.017$ \\
\bottomrule
\end{tabular}
\end{table}


\subsection{Zeros}\label{sec:zeros}

\begin{proposition}\label{prop:zerofree}
\begin{enumerate}
\item There is $\sigma_1$ such that $\zeta_F(s) \ne 0$ for $\operatorname{Re}s \ge \sigma_1$.
\item For every $\varepsilon > 0$ and $\sigma_2 < -n_0-\varepsilon$ there is $t_0$ such that $\zeta_F(s) \ne 0$ for
$\sigma_2 \le \operatorname{Re}s \le -n_0-\varepsilon$, $|\operatorname{Im}s| \ge t_0$.
\end{enumerate}
\end{proposition}

\begin{proof}
(1) If $k_1$ is the least $k$ with $F(k) \ne 0$, then $k_1^s\zeta_F(s) \to F(k_1)$ as $\operatorname{Re}s \to \infty$,
uniformly in $\operatorname{Im}s$. (2) In the strip, write $\zeta_F$ as in the proof of
Proposition~\ref{prop:growth}. By the functional equation and Stirling's formula,
$|\zeta(s+n_0)| \asymp t^{1/2-\sigma-n_0}|\zeta(1-s-n_0)| \gg t^{1/2+\varepsilon}$, since
$\operatorname{Re}(1-s-n_0) \ge 1+\varepsilon$. Each $\zeta(s+n)$ with $n > n_0$ is
$O(t^{1/2-\sigma-n_0-1}+t^{\varepsilon/2})$, and everything else is bounded, so $\zeta_F(s) \ne 0$ for large $t$.
\end{proof}

\begin{theorem}\label{thm:count}
For $0 < \varepsilon < 1$ let $N_F(T)$ be the number of zeros $\rho$ of $\zeta_F$, counted with multiplicity, with
$\operatorname{Re}\rho \ge -n_0-\varepsilon$ and $0 < \operatorname{Im}\rho \le T$. Then
\[
  N_F(T) = \frac{T}{2\pi}\log\frac{T}{2\pi e k_1} + O(\log T),
\]
where $k_1$ is the least $k$ with $F(k) \ne 0$. In particular, if $F(1) \ne 0$, then $N_F(T) - N(T) = O(\log T)$,
where $N(T)$ counts the zeros of $\zeta$.
\end{theorem}

\begin{proof}[Proof sketch]
By Proposition~\ref{prop:zerofree} the zeros counted lie, apart from finitely many, in the rectangle
$-n_0-\varepsilon \le \sigma \le \sigma_1$, $t_0 \le t \le T$. Apply the argument principle to its boundary. On the
right side the change of argument is $-T\log k_1 + O(1)$, since $\operatorname{Re}\bigl(k_1^s\zeta_F(s)/F(k_1)\bigr) > \frac12$
there. On the
left side, by the proof of Proposition~\ref{prop:zerofree}, $\zeta_F(s) = a_{n_0}\zeta(s+n_0)(1 + o(1))$, so the change
of argument is that of $\zeta(-\varepsilon+it)$ up to $O(1)$, which by the functional equation is that of
$\chi(-\varepsilon+it)$ up to $O(1)$, namely $-T\log(T/2\pi e) + O(1)$. On the top side the change of argument is
$O(\log T)$ by the usual Jensen argument \cite[\S9.4]{Tit86}, using the polynomial growth of $\zeta_F$
(Proposition~\ref{prop:growth}) and $|\zeta_F(\sigma_1+iT)| \gg 1$.
\end{proof}

(For $\zeta(s) - 1$, i.e.\ $k_1 = 2$, this is Landau's count of the $1$-points of $\zeta$.) So if $F(1) \ne 0$,
$\zeta_F$ has as many zeros as $\zeta$, up to $O(\log T)$. Where are they? On the line
$\operatorname{Re}s = \frac12 - n_0$ we have $\zeta_F(s) = a_{n_0}\zeta(s+n_0) + B(s)$, where
$B(s) = \zeta_F(s) - a_{n_0}\zeta(s+n_0,k_0) - \sum_{k=1}^{k_0-1}F(k)k^{-s}$ is an absolutely convergent Dirichlet series
for $\operatorname{Re}s > -n_0$, with coefficients $O(k^{-n_0-1})$. The zeros of $\zeta_F$ are thus the solutions of
$\zeta(w) = -B(w-n_0)/a_{n_0}$: ``$a$-points'' of $\zeta$ with a variable, bounded, almost periodic $a$. Levinson
\cite{Lev75} proved that for constant $a$ all but $O(N(T)/\log\log T)$ of the $a$-points with
$0 < \operatorname{Im}s \le T$ lie in $|\operatorname{Re}s - \frac12| < \varepsilon$. This suggests:

\begin{conjecture}\label{conj:cluster}
For every $\varepsilon > 0$, all but $o(T\log T)$ of the zeros of $\zeta_F$ with $0 < \operatorname{Im}\rho \le T$ and
$\operatorname{Re}\rho \ge -n_0 - \varepsilon$ satisfy $|\operatorname{Re}\rho - (\frac12 - n_0)| < \varepsilon$.
\end{conjecture}

There is no reason to expect a Riemann hypothesis: without an Euler product, and as for the Davenport--Heilbronn
examples \cite{DH36}, there should be $\gg T$ zeros off the line.

\subsection{Numerical experiments}
For Khinchin's $Z$ ($n_0 = 2$), a search in $-6 \le \sigma \le 3$, $0 < t \le 50$ finds exactly one zero near each
zero $\frac12 + i\gamma_n$ of $\zeta(s+2)$, i.e.\ near $-\frac32 + i\gamma_n$, but displaced to the right by about
$0.5$--$0.7$ (Table~\ref{tab:zeros}), and three zeros of small height. To first order, the zero near
$-\frac32 + i\gamma_n$ is displaced by $-B(-\frac32 + i\gamma_n)/\zeta'(\frac12+i\gamma_n)$, which is small only when
$|\zeta'|$ is large; at these heights the first-order prediction gives the direction but underestimates the shift.
Higher up the displacement decreases slowly, in line with Conjecture~\ref{conj:cluster}: near
$\gamma_{100} = 236.52$, $\gamma_{1000} = 1419.42$ and $\gamma_{3000} = 3533.33$ the zeros of $Z$ have real parts
$-1.203$, $-1.138$ and $-1.273$.

\begin{table}[t]
\centering\small
\caption{The zeros $\rho$ of Khinchin's $Z(s)$ with $5 < \operatorname{Im}\rho \le 50$, next to the ordinates
$\gamma_n$ of the zeros of $\zeta$.}
\label{tab:zeros}
\begin{tabular}{rrr@{\qquad}rrr}
\toprule
$\gamma_n$ & $\operatorname{Re}\rho$ & $\operatorname{Im}\rho$ & $\gamma_n$ & $\operatorname{Re}\rho$ & $\operatorname{Im}\rho$ \\
\midrule
$14.1347$ & $-0.80962$ & $14.10628$ & $37.5862$ & $-0.97200$ & $37.62163$ \\
$21.0220$ & $-0.87839$ & $21.04380$ & $40.9187$ & $-0.97864$ & $40.96155$ \\
$25.0109$ & $-0.90838$ & $24.96595$ & $43.3271$ & $-0.99413$ & $43.22684$ \\
$30.4249$ & $-0.93768$ & $30.50982$ & $48.0052$ & $-1.00783$ & $48.16616$ \\
$32.9351$ & $-0.95140$ & $32.84059$ & $49.7738$ & $-1.00367$ & $49.65618$ \\
\bottomrule
\end{tabular}
\end{table}

The three zeros of small height, $-1.5119 + 1.6018i$, $-3.0436 + 2.1068i$ and $-4.5285 + 2.4371i$, start a chain
continuing to the left between the poles at the negative integers, with real parts spaced by about $1.3$--$1.5$:
approximately $-5.9950 + 2.6727i$, $-7.4513 + 2.8396i$, \dots, $-11.7643 + 3.0414i$, where the imaginary part is
largest, and then $-14.5900 + 2.9928i$, $-20.1220 + 2.5933i$, $-26.8343 + 1.6953i$, $-30.7659 + 0.9964i$,
$-33.3916 + 0.4946i$. After $24$ complex zeros in the upper half-plane (and their conjugates) the chain reaches the real axis, where it
becomes the pair of real zeros $-34.7863$ and $-34.4495$. Further left, a scan of the real axis down to $-70$ (with step $\frac1{40}$) finds
one zero in each interval $(-2j, -2j+1)$ with $38 \le 2j \le 70$, close to the even integer, and no others: $-37.7571$, $-39.7968$, $-41.8033$, \dots,
$-67.8796$, $-69.8831$, with distance to $-2j$ slowly decreasing from $0.24$ to $0.12$. These play the role of the
trivial zeros $-2j$ of $\zeta(s+2)$. So, numerically, $Z$ has only finitely many non-real zeros of small height, and
its remaining trivial zeros are real. Heuristically, far to the left the first dual term of Poisson summation
dominates, and the zeros approach those of $s \mapsto \widehat G_s(1)$, a Mellin transform of $F(x)\cos(2\pi x)$; for
$F = 1$ these are the zeros of $\sin(\pi s/2)$. Since we cannot exclude other functions $\zeta_F$ with infinitely
many non-real zeros of bounded height, Theorem~\ref{thm:count} restricts to $\operatorname{Re}\rho \ge -n_0-\varepsilon$.

\section{Questions}
\begin{enumerate}
\item \emph{Functional equations.} For which $F$ is $\widehat G_s$ again of the form $|\xi|^{s'}F'(|\xi|)$ with
$F'$ regular, giving a functional equation relating $\zeta_F$ and $\zeta_{F'}$? Besides powers, the inhomogeneous
Epstein (Chowla--Selberg) sums $\sum_{k=1}^{\infty}(k^2+a^2)^{-s}$, with modified Bessel functions on the dual side
\cite{CS67,BGMWZ13}, are of this kind, but only as diagonals of two-variable families.
\item \emph{Zeros.} Prove Conjecture~\ref{conj:cluster}, and determine whether a positive proportion of the zeros
lie off the line $\operatorname{Re}s = \frac12-n_0$. Does every $\zeta_F$ have only finitely many non-real zeros with $\operatorname{Re}\rho < -n_0-\varepsilon$, as
Khinchin's $Z$ appears to have?
\item \emph{Value distribution.} Is $\zeta_F$ universal (in the sense of Voronin) in the strip
$\frac12 - n_0 < \operatorname{Re}s < 1 - n_0$? Is it hypertranscendental?
\end{enumerate}

\begin{thebibliography}{99}
\bibitem{BGMWZ13} J.~M.~Borwein, M.~L.~Glasser, R.~C.~McPhedran, J.~G.~Wan and I.~J.~Zucker, \emph{Lattice Sums Then and Now}, Cambridge University Press, 2013.
\bibitem{CS67} S.~Chowla and A.~Selberg, On Epstein's zeta-function, J.\ Reine Angew.\ Math.\ 227 (1967), 86--110.
\bibitem{DH36} H.~Davenport and H.~Heilbronn, On the zeros of certain Dirichlet series, J.\ London Math.\ Soc.\ 11 (1936), 181--185 and 307--312.
\bibitem{FGD95} P.~Flajolet, X.~Gourdon and P.~Dumas, Mellin transforms and asymptotics: harmonic sums, Theoret.\ Comput.\ Sci.\ 144 (1995), 3--58.
\bibitem{FV96} P.~Flajolet and I.~Vardi, Zeta function expansions of classical constants, manuscript, 1996.
\bibitem{HL23} G.~H.~Hardy and J.~E.~Littlewood, The approximate functional equation in the theory of the zeta-function, with applications to the divisor-problems of Dirichlet and Piltz, Proc.\ London Math.\ Soc.\ (2) 21 (1923), 39--74.
\bibitem{Joh26} F.~Johansson, Computing infinite sums of holonomic functions, preprint, 2026.
\bibitem{Kat97} M.~Katsurada, An application of Mellin--Barnes' type integrals to the mean square of Lerch zeta-functions, Collect.\ Math.\ 48 (1997), 137--153.
\bibitem{Khi35} A.~Khintchine, Metrische Kettenbruchprobleme, Compositio Math.\ 1 (1935), 361--382.
\bibitem{Khi64} A.~Ya.~Khinchin, \emph{Continued Fractions}, University of Chicago Press, 1964.
\bibitem{Lev75} N.~Levinson, Almost all roots of $\zeta(s) = a$ are arbitrarily close to $\sigma = 1/2$, Proc.\ Nat.\ Acad.\ Sci.\ USA 72 (1975), 1322--1324.
\bibitem{Ost20} A.~Ostrowski, \"Uber Dirichletsche Reihen und algebraische Differentialgleichungen, Math.\ Z.\ 8 (1920), 241--298.
\bibitem{Rie59} B.~Riemann, Ueber die Anzahl der Primzahlen unter einer gegebenen Gr\"osse, Monatsber.\ Berliner Akad., 1859.
\bibitem{Rub05} M.~Rubinstein, Computational methods and experiments in analytic number theory, in \emph{Recent Perspectives in Random Matrix Theory and Number Theory}, LMS Lecture Note Ser.\ 322, Cambridge University Press, 2005, 425--506.
\bibitem{Ste07} J.~Steuding, \emph{Value-Distribution of $L$-Functions}, Lecture Notes in Math.\ 1877, Springer, 2007.
\bibitem{Tit86} E.~C.~Titchmarsh, \emph{The Theory of the Riemann Zeta-Function}, 2nd ed., revised by D.~R.~Heath-Brown, Oxford University Press, 1986.
\end{thebibliography}

\end{document}
